\section{Introducción}

El objetivo de este proyecto es diseñar un sistema fotovoltaico aislado que mantenga en funcionamiento las cargas críticas de la vivienda de Abner López, en Santa Rosa de Oreamuno, cuando falla el servicio eléctrico. La vivienda tiene servicio de JASEC, por lo que el sistema se planteó como respaldo: las cargas críticas se alimentan desde un subtablero propio, normalmente con energía solar y el banco de baterías, y la red queda como apoyo. El banco se dimensionó para dos días de autonomía.

El diseño sigue el procedimiento de sistemas aislados del curso: caracterización del consumo, recurso solar, balance de energía, banco de baterías, generador fotovoltaico, regulador de carga e inversor. Parte de los datos de placa de los equipos de la casa, de la factura de JASEC, de la base de datos PVGIS y de las fichas técnicas de los equipos seleccionados.
